对稳恒电流，安培环路定理写成
 \(\oint_L \vb*{H}\cdot d\vb*{l}=I_c\)
足以描述这一情形。但若考虑给平行板电容器充电，连接导线中有传导电流 \(I_c\)，极板间却没有真实载流子穿越空隙。若仍只把右边写成传导电流，就会出现“同一条环路用不同跨面时结果不一致”的矛盾。

这时的关键不是“再硬加一项试试”，而是先考察电容器板间实际随时间变化的量：虽然没有电子穿过空隙，但极板上的自由电荷 \(q(t)\) 仍随时间增加，板间电场和电位移 \(\vb*{D}\) 也随之变化。若以平行板电容器为例，忽略边缘效应时
 \(D=\frac{q}{S} \implies \Phi_D=DS=q.\)
于是
 \(\dv{\Phi_D}{t}=\dv{q}{t}=I_c.\)
这说明极板间虽然没有传导电流，却有一个与外电路电流严格对应的“通量变化率”。麦克斯韦引入位移电流，正是将这一连续变化纳入磁场的源。

麦克斯韦为解决这个矛盾引入\textbf{位移电流}：
 \(I_d=\dv{\Phi_D}{t}, \qquad \vb*{j}_d=\pdv{\vb*{D}}{t}.\)
于是安培定律改写为
 \(\oint_L \vb*{H}\cdot d\vb*{l}=I_c+I_d.\)
若把传导电流和位移电流合在一起，就得到\textbf{全电流}：
 \(\vb*{j}_{\mathrm{all}}=\vb*{j}_c+\vb*{j}_d =\vb*{j}_c+\pdv{\vb*{D}}{t},\)
 \(I_{\mathrm{all}}=I_c+I_d =I_c+\dv{\Phi_D}{t}.\)
麦克斯韦修正的关键不只是“多补一项”，而是使同一条环路取不同跨面时，总电流始终一致。也正因为如此，非稳恒情形下电流的连续性才不会被电容器这种结构打断。

\begin{note}
判断位移电流是否存在，不看“中间有没有导线”，而看\textbf{电位移通量是否随时间变化}。平行板电容器中的两种情况最能说明这一判据：
\begin{enumerate}
 \item 若电容器充完电后断开电源，再缓慢拉大板距，则自由电荷 \(Q\) 不变，忽略边缘效应时 \(D=Q/S\) 也不变，所以 \(\pdv{\vb*{D}}{t}=0\)，极板间\textbf{没有}位移电流。
 \item 若电容器始终接在电源上，再拉大板距，则电压 \(U\) 被电源维持不变，而 \(E\approx U/d\) 会随 \(d\) 改变，从而 \(D=\varepsilon E\) 也改变，因此 \(\pdv{\vb*{D}}{t}\neq 0\)，极板间\textbf{存在}位移电流。
\end{enumerate}
所以，“有没有位移电流”的判断取决于\textbf{谁被保持不变、谁随时间变化}，而不是“有没有真实载流子穿过去”。
\end{note}
